\section{Runtime Analysis}
\label{sec:runtimes}

Here, we provide a detailed asymptotic analysis of runtimes for different tensor products. We consider 3 different settings.
\begin{itemize}
    \item 
    \textbf{Single Input, Single Output (SISO)}:
    
    Here we are computing only one path $[\ell_1,\ell_2,\ell_3]$ where $\ell_i\in\O(L).$
    \[\ell_1\times\ell_2\to\ell_3\] 
    \item \textbf{Single Input, Multiple Output (SIMO)}:
    
    Here we fix $\ell_1,\ell_3$ but allow all possible irreps generated by the respective tensor products.
    \[\ell_1 \times \ell_2 \to Z\]
    \item \textbf{Multiple Input, Multiple Output (MIMO)}: 
    
    Here we only bound the $L$ that the tensor products use but allow full capacity for the input and output irreps. In the case of CGTP, we can have an arbitrary number of copies of each irrep but we assume we only use single copies of each irrep in the input.
    \[Z\times W\to Z\]
\end{itemize}
In the SISO and SIMO settings, the asymptotic runtimes of different tensor products are directly comparable. However, in the MIMO setting, we lose expressivity in some tensor products. This is discussed more in \autoref{sec:expressivity}. Note the MIMO setting is what one would typically want to use in practice. Hence the runtimes reported in \autoref{tab:asymptotic_runtimes_IO} are for the MIMO setting.

\subsection{Clebsch-Gordan Tensor Product}
\label{sec:CGTP_runtime}
The tensor product operation is defined as:
\begin{align}
\label{eqn:tp} (\xl{1} \cgtp \xl{2})^{(\ell_3)}_{m_3} =\sum_{m_1=-l_1}^{l_1}\sum_{m_2=-l_2}^{l_2}
C^{(\ell_3,m_3)}_{\ell_1,m_1,\ell_2,m_2}\xl{1}_{m_1}\xl{2}_{m_2}
\end{align}


where $C$ denotes the Clebsch-Gordan (CG) coefficients which can be precomputed.

\subsubsection{Naive runtime}
Let $L=\max(\ell_1,\ell_2,\ell_3)$.
From \autoref{eqn:tp}, for each $m_3$, we would need to sum over $m_1,m_2$ which range from $-\ell_1$ to $\ell_1$ and $-\ell_2$ to $\ell_2$ respectively. Hence, we expect $\O(L^2)$ operations. To compute the values for all $m$ which range from $-\ell_3$ to $\ell_3$, we see that computing a single $\ell_1\otimes\ell_2\to\ell_3$ tensor product requires $\O(L^3)$ operations.

\subsubsection{Optimized runtime with sparsity}
However, the CG coefficients are sparse. In the complex basis for the irreps, $C^{(\ell_3,m)}_{\ell_1,m_1,\ell_2,m_2}$ is nonzero only if $m_1+m_2=m_3$. Transforming to the real basis for the irreps, this condition becomes $\pm m_1\pm m_2=m_3$. In either case for a fixed $m_1$ and $m_3$, we only ever need to sum over a constant number of $m_2$'s rather than $\O(L)$ of them as naively expected. Therefore an implementation taking this sparsity into account gives us a runtime of $\O(L^2)$. This optimization was noted in \citet{cobbefficient}.


\subsection{Gaunt Tensor Product}
\label{sec:GTP_runtimes}
The Gaunt Tensor Product (GTP) is based on the decomposition of a product of spherical harmonic functions back into spherical harmonics \cite{gaunt}. In particular, suppose one of our inputs $\xl{1}$ transforms as a direct sum of irreps up to some cutoff $L$ (ie. $\ell_1$ ranges from $0,\ldots,L$). We can view these irreps as coefficients of spherical harmonics which gives a spherical signal $F_1(\theta,\varphi)=\sum_{\ell_1,m_1}\xl{1}_{m_1}Y_{\ell_1,m_1}(\theta,\varphi)$. We similarly construct $F_2(\theta,\varphi)=\sum_{\ell_2,m_2}\xl{2}_{m_2}Y_{\ell_2,m_2}(\theta,\varphi)$.

Taking the product of these spherical signals gives a new signal $F_3(\theta,\varphi)=F_1(\theta,\varphi)F_2(\theta,\varphi)$. This new signal can be decomposed into spherical harmonics which we use to define the GTP. This results in
\begin{align}
F_3(\theta,\varphi)=\sum_{\ell_3,m_3}(\xl{1} \gtp \xl{2})^{(\ell_3)}_{m_3}Y_{\ell_3,m_3}(\theta,\varphi).\end{align}



\subsubsection{2D Fourier basis}
\citet{gaunt} describe an implementation which decomposes spherical harmonics into a 2D Fourier basis in their original paper introducing GTP. This also turns out to be the same implementation in \citet{xin2021fast}. We describe their procedure here.

Note that for any $\ell\leq L$ we can always write the spherical harmonics in the 2D Fourier basis:
\begin{align}
    Y_{\ell,m}(\theta,\varphi)=\sum_{-L\leq u,v\leq L}y^{\ell,m}_{u,v}e^{i(u\theta+v\varphi)}
\end{align}
for some coefficients $y^{\ell,m}_{u,v}$.

Hence, any signal $\mathbf{x}^{(\ell)}_m$ can be encoded as
\begin{align}
F_1(\theta,\varphi)=\sum_{\ell=0}^L\sum_{m=-\ell}^\ell\sum_{-L\leq u,v\leq L}\mathbf{x}^{(\ell)}_m y^{\ell,m}_{u,v}e^{i(u\theta+v\varphi)}=\sum_{-L\leq u,v\leq L}\left(\sum_{\ell=0}^L\sum_{m=-\ell}^\ell\mathbf{x}^{(\ell)}_m y^{\ell,m}_{u,v}\right)e^{i(u\theta+v\varphi)}.
\end{align}
We identify the encoding
\begin{align}
\mathbf{x}_{u,v}=\sum_{\ell=0}^L\sum_{m=-\ell}^\ell\mathbf{x}^{(\ell)}_m y^{\ell,m}_{u,v}. 
\end{align}
One can observe that the $y^{\ell,m}_{u,v}$ are sparse and only nonzero when $m=\pm v$. Therefore, finding $\mathbf{x}_{u,v}$ if we have a set of irreps is $\O(L)$ and it is $O(1)$ if we only want one irrep. Because there are $\O(L^2)$ possible values for $u,v$, encoding into the 2D Fourier is $\O(L^3)$ if we encode all irreps up to $L$ or $\O(L^2)$ if encoding a single irrep.

For 2 functions of $\theta,\varphi$ encoded using a 2D Fourier basis $\mathbf{x}^1_{u,v},\mathbf{x}^2_{u,v}$, we can compute their product using a standard 2D FFT in $\O(L^2\log L)$ time. This gives some output encoded as $\mathbf{y}_{u,v}$ where now $u,v$ range from $-2L,\ldots,2L$ to capture all information.

Finally, we decode the resulting function in the 2D Fourier basis back into a spherical harmonic basis to extract the output irreps. Suppose $-L\leq u,v\leq L$. We can always write
\begin{align}
e^{i(u\theta+v\varphi)}=F^\perp_{u,v}(\theta,\varphi)+\sum_{\ell=0}^L\sum_{m=-\ell}^\ell z^{\ell,m}_{u,v}Y_{\ell,m}(\theta,\varphi)
\end{align}
where $F^\perp_{u,v}(\theta,\varphi)$ is some function in the space orthogonal to that spanned by the spherical harmonics. By construction, our output signal is always in the space spanned by the spherical harmonics so the orthogonal parts cancel. Hence we can write
\begin{align}
    \sum_{-2L\leq u,v\leq 2L}\mathbf{y}_{u,v}e^{i(u\theta+v\varphi)}= & \sum_{-2L\leq u,v\leq 2L}\mathbf{y}_{u,v}\sum_{\ell=0}^L\sum_{m=-\ell}^\ell z^{\ell,m}_{u,v}Y_{\ell,m}(\theta,\varphi)\\
    = & \sum_{\ell=0}^L\sum_{m=-\ell}^\ell \left(\sum_{-2L\leq u,v\leq 2L}\mathbf{y}_{u,v} z^{\ell,m}_{u,v}\right)Y_{\ell,m}(\theta,\varphi)
\end{align}
Hence we identify:
\begin{align}
\mathbf{y}^\ell_m=\sum_{-2L\leq u,v\leq 2L}\mathbf{y}_{u,v} z^{\ell,m}_{u,v}.
\end{align}
Once again, we can note that $z^{\ell,m}_{u,v}$ must be sparse and is only nonzero when $v=\pm m$. Hence, evaluating the above takes $\O(L)$ time since we sum over $\O(L)$ values of $u$ paired with constant number of $v$'s. If we only extract one irrep, then we range over $\O(L)$ values of $m$ giving $\O(L^2)$ runtime. If we extract all irreps up to $2L$ this becomes $\O(L^3)$.

\subsubsection{Grid tensor product}
\label{sec:grid_runtime}
Rather than use a 2D Fourier basis, we can instead represent the signal by directly giving its value for a set of points on the sphere. Quadrature on the sphere is a well-studied topic \citep{beentjes2015quadrature,lebedev1976quadratures}; in general, $\O(L^2)$ points are needed to exactly integrate spherical harmonics upto degree $L$ \citep{mclaren1963optimal}. For this section, consider a product grid on the sphere formed by the Cartesian product of two 1D grids for $\theta$ and $\varphi$ with $\O(L)$ points each, for a total of $\O(L^2)$ points. 

We can write:
\begin{align}
F_1(\theta_{j},\varphi_{k})=\sum_{\ell=0}^L\sum_{m=-\ell}^\ell \mathbf{x}^{(\ell)}_m Y_{\ell,m}(\theta_j,\varphi_k)=\sum_{\ell=0}^L\sum_{m=-\ell}^\ell \mathbf{x}^{(\ell)}_m N_{\ell,m}P^{m}_\ell(\cos(\theta_j))cs_m(\varphi_k)
\end{align}
where $N_{\ell,m}$ is some normalization factor, $P_\ell^m$ are the associated Legendre polynomials, and
\begin{align}
    cs_m(\varphi)=\begin{cases}
    \sin(|m|\varphi) & m<0\\
    1 & m=0\\
    \cos(m\varphi) & m>0
\end{cases}.
\end{align}
We note that we can first evaluate
\begin{align}
   g_{m}(\theta_j)=\sum_{\ell=0}^L\mathbf{x}^{(\ell)}_m N_{\ell,m}P^{m}_\ell(\cos(\theta_j))
\end{align}
where we set $P^m_\ell=0$ if $m>\ell$. If we have a set of irreps up to $L$ then we do the summation and this takes $\O(L)$ time. If we only have one irrep to encode then this takes $O(1)$ time. But we also have $\O(L)$ values of $\theta_j$ on the grid and $\O(L)$ values of $m$ to evaluate. This gives $\O(L^3)$ runtime to encode onto the grid for irreps up to $L$ and $\O(L^2)$ for a single irrep. Finally evaluating 
\begin{align}
F_1(\theta_j,\varphi_k)=\sum_{m=-\ell}^\ell g_m(\theta_j) cs_m(\varphi_k)
\end{align}
for a set of $\varphi_k$ can be done through a FFT in $\O(L\log L)$ time for each $\theta_j$ giving $\O(L^2\log L)$ total. Hence we see encoding onto the sphere takes $\O(L^3)$ time for irreps up to $L$ and $\O(L^2\log L)$ time for a single irrep.

For the multiplication of signals, we just have elementwise multiplication $F_3(\theta_k,\varphi_k)=F_1(\theta_k,\varphi_k)\cdot F_2(\theta_k,\varphi_k)$. Since there are $\O(L^2)$ grid points this takes $\O(L^2)$ time.

Finally, we decode the signal back into irreps. To do so we use the fact that
\begin{align}\mathbf{f}^{(\ell)}_m=\sum_{j,k}a_j F(\theta_j,\varphi_k)Y_{\ell,m}(\theta_j,\varphi_k)\end{align}
for some coefficients $a_j$. This is essentially performing numerical integration of our signal against a spherical harmonic. Once again using the factorization of the spherical harmonics we get
\begin{align}\mathbf{f}^{(\ell)}_m=\sum_j\left(\sum_{k} F(\theta_j,\varphi_k)cs_m(\varphi_k)\right)a_j N_{\ell,m}P^{m}_\ell(\cos(\theta_j)).\end{align}
The inner sum in parentheses can be computed in $\O(L)$ time and we need to compute it for $\O(L^2)$ values of $\theta_j,m$ pairs giving a runtime of $\O(L^3)$. Of course, we note that $cs$ really is just sines and cosines so alternatively we can use FFT which takes $\O(L^2\log L)$ total. Computing the outer sum takes $\O(L)$ since we sum over $\O(L)$ values of $j$. For a single irrep there are $\O(L)$ values of $j$ giving $\O(L^2)$ for the outer sum. For irreps up to $\ell$ there are $\O(L^2)$ pairs of $\ell,m$ giving $\O(L^3)$ runtime for the outer sum. In total, we see going from the grid to the coefficients takes $\O(L^2\log L)$ for a single irrep and $\O(L^3)$ for all irreps.

However, it turns out that the associated Legendre polynomials have recurrence properties which can be exploited to make transforming a set of irreps up to $L$ to the grid and a set of irreps up to $L$ back from the grid asymptotically more efficient \cite{healy2003ffts}. The runtime for this algorithm which we will call S2FFT is $\O(L^2\log^2 L)$.

\subsection{Matrix Tensor Product}
\label{sec:FTP_runtime}

Here we describe and analyze the time complexity of matrix tensor product. Let $L_1,L_2$ be the max $\ell$'s of the inputs and $L_3$ be the max $\ell$ of the outputs. We pick some $\tilde{\ell}=\lceil\max(L_1,L_2,L_3)/2\rceil$ so that $\tilde{\ell}\otimes\tilde{\ell}$ when decomposed into irreps can contain all irreps of the inputs and outputs. Note in principle we could always choose larger $\tilde{\ell}$.

In the following runtime analysis, we assume $L_1=L_2=L$, $\tilde{l}=L$, and $L_3=2L$.

\subsubsection{Naive runtime}
The first step of MTP is to convert our input irreps into a tensor product rep using Clebsch-Gordan coefficients as
\begin{align}
    \mathbf{X}^{(\ell)}_{m_1,m_2} = \sum_{m_3=-\ell}^{\ell} C^{\ell_3,m_3}_{\tilde{\ell},m_1,\tilde{\ell},m_2}\xl{}_{m_3}\\
    \mathbf{Y}^{(\ell)}_{m_1,m_2} = \sum_{m_3=-\ell}^{\ell} C^{\ell_3,m_3}_{\tilde{\ell},m_1,\tilde{\ell},m_2}\yl{}_{m_3}.
\end{align}
Naively we sum over $\O(L)$ values of $m_3$ and need to do the computation for $\O(L^2)$ possible pairs of $m_1,m_2$. This gives $\O(L^3)$ naive runtime for converting a single irrep into a tensor product rep. To do so for all irreps up to $L$ the takes $\O(L^4)$ time.

We can then sum over tensor product reps to create 
\begin{align}
    \mathbf{X}=\sum_{\ell}\mathbf{X}^{(\ell)}\qquad
    \mathbf{Y}=\sum_{\ell}\mathbf{Y}^{(\ell)}.
\end{align}
There are $\O(L)$ matrices to sum over if we have irreps up to $L$. Summing matrices takes $\O(L^2)$ time since our matrices are size $\O(L)\times \O(L)$. Hence, this takes $\O(L^3)$ time if we have irreps up to $L$. If we have a single irrep then we do not need to do anything.

We then multiply the matrices giving $\mathbf{Z}=\mathbf{X}\mathbf{Y}$. Using the naive matrix multiplication algorithm requires $\O(L^3)$ runtime.

Finally we can use Clebsch-Gordan to extract individual irreps giving
\begin{align}
    (\x\ftp\y)^{(\ell_3)}_{m_3}=\sum_{m_1=-\ell_1}^{\ell_1}\sum_{m_2=-\ell_2}^{\ell_2} C^{(\ell_3,m_3)}_{\ell_1,m_1,\ell_2,m_2}\mathbf{Z}_{m_1,m_2}.
\end{align}
Again, naively we sum over $\O(L^2)$ pairs of $m_1,m_2$ and need to evaluate $\O(L)$ values of $m_3$ for $\O(L^3)$ conversion for single irrep. If we want all irreps up to $2L$ then we need $\O(L^4)$.

\subsubsection{Optimized runtime with sparsity}
Similar to the CGTP, we can take sparsity of the Clebsch-Gordan coefficients into account. We have nonzero values only if $\pm m_1\pm m_2=m_3$. Hence in the encoding step, for fixed $m_1,m_2$ we only need to sum over constant number of $m_3$ instead of $\O(L)$. This gives a reduction of $L$ in encoding to tensor product rep. Similarly in the decoding step, we see for fixed $m_3$ we only need to sum over $\O(L)$ pairs of $m_1,m_2$. This gives a reduction of $L$ as well in decoding back into irreps.

\subsection{Asymptotic runtimes in different settings}

\begin{table}[h]
\caption{Asymptotic runtimes of various tensor products for different output settings. The best performing tensor products for each output settings are highlighted {\color{highlight} in green}. In the MIMO setting, the Clebsch-Gordan tensor products are highlighted {\color{warning} in red} to indicate that they can output irreps with multiplicity $ > 1$ , unlike the Gaunt tensor products.}
\vspace{0.3cm}
\label{tab:asymptotic_runtimes}
    \centering    
    \scalebox{0.7}{\begin{tabular}{cccc}
        \toprule
        Tensor Product & SISO & SIMO & MIMO \\
         \midrule
        Clebsch-Gordan (Naive) & $\O(L^3)$ & $\O(L^4)$ & $\color{warning} \O(L^6)$ \\
        Clebsch-Gordan (Sparse) & $\color{highlight} \O(L^2)$ & $\O(L^3)$ & $\color{warning} \O(L^5)$ \\
        Gaunt (Original) & $\O(L^2\log L)$ & $\O(L^3)$ & $\O(L^3)$ \\
        Gaunt (Naive Grid) & $\O(L^2\log L)$ & $\O(L^3)$ & $\O(L^3)$ \\
        Gaunt (S2FFT Grid) & $\O(L^2\log L)$ & $\color{highlight} \O(L^2\log^2 L)$ & $\color{highlight} \O(L^2 \log^2 L)$ \\
        Matrix (Naive)& $\O(L^3)$ & $\O(L^4)$ & $\O(L^4)$ \\
        Matrix (Sparse)& $\O(L^3)$ & $\O(L^3)$ & $\O(L^3)$ \\
        \bottomrule
    \end{tabular}}
\end{table}